Model-based reinforcement learning (RL) turns experience into a model and then uses the model to generate extra updates. In Dyna~\cite{sutton1991dyna} each real transition is followed by $n$ simulated updates drawn from the learned model, so $n$ is a dial between pure model-free learning (Q-learning~\cite{watkins1992q}, $n=0$) and heavy planning. The benefit is sample efficiency; the cost is that every planning update trusts the model. When the environment changes after the model was learned (a \emph{stale} model) or the transitions are stochastic while the model stores a single outcome (a \emph{wrong} model), more planning amplifies the error. The same tension drives modern deep model-based RL, where short rollouts~\cite{janner2019when} and probabilistic ensembles~\cite{chua2018deep} are used to limit compounding model error~\cite{moerland2020model}, and where benchmarking studies find large differences across algorithms and tasks~\cite{wang2019benchmarking}.

The tabular case is the cleanest place to isolate the effect, and the blocking and shortcut mazes of Sutton's Dyna work motivate Dyna-Q+~\cite{sutton1990integrated} (an exploration bonus for transitions not tried recently) as a remedy. Textbook comparisons are, however, illustrative curves: they do not sweep the number of planning steps, do not include stochastic dynamics, and do not report variability over seeds. This paper runs that controlled experiment on CPU. We ask how $n$ trades early and post-change reward against model error for four reimplemented agents, and state five falsifiable hypotheses (Sec.~\ref{sec:setup}). Our contributions are an honest, small-scale measurement of the $n$-dependence on five maze regimes with 10--20 seeds, an ablation and sensitivity study of the two Dyna-Q+ ingredients, and the report of what did not replicate: no staleness-driven decline of Dyna-Q's post-change reward with $n$ appears, and Dyna-Q+ is brittle to its bonus weight.
